\documentclass{article}
\usepackage{amsmath,amssymb,xurl}
\usepackage[hidelinks]{hyperref}
% Shared commands for notes in docs/paper-gaps/.

\DeclareUnicodeCharacter{2080}{\ifmmode{}_{0}\else\textsubscript{0}\fi}
\DeclareUnicodeCharacter{2081}{\ifmmode{}_{1}\else\textsubscript{1}\fi}
\DeclareUnicodeCharacter{2082}{\ifmmode{}_{2}\else\textsubscript{2}\fi}
\DeclareUnicodeCharacter{2099}{\ifmmode{}_{n}\else\textsubscript{n}\fi}

\newcommand{\Lean}{\textsc{Lean}}
\ExplSyntaxOn
\NewDocumentCommand{\leanid}{m}{
  \ifmmode
    \textsf{\detokenize{#1}}
  \else
    \group_begin:
    \urlstyle{sf}
    \tl_set:Nn \l_tmpa_tl {#1}
    \tl_replace_all:Nnn \l_tmpa_tl {\_} {_}
    \exp_args:NV \path \l_tmpa_tl
    \group_end:
  \fi
}
\ExplSyntaxOff
\providecommand{\tr}{\operatorname{tr}}
\newcommand{\tnleanIssueBase}{https://github.com/LionSR/TNLean/issues/}
\newcommand{\tnleanPrBase}{https://github.com/LionSR/TNLean/pull/}
\newcommand{\tnleanDocBase}{https://sirui-lu.com/TNLean/docs/}
\newcommand{\ghissue}[1]{\href{\tnleanIssueBase#1}{GitHub issue \##1}}
\newcommand{\ghpr}[1]{\href{\tnleanPrBase#1}{PR \##1}}
\newcommand{\doclink}[1]{\href{\tnleanDocBase#1.html}{\path{#1}}}

% Dirac notation and the identity operator, as in blueprint/src/macros/common.tex.
% \providecommand keeps note-local definitions authoritative.
\usepackage{dsfont}
\providecommand{\ket}[1]{|#1\rangle}
\providecommand{\bra}[1]{\langle #1|}
\providecommand{\braket}[2]{\langle #1 | #2 \rangle}
\providecommand{\ketbra}[2]{|#1\rangle\!\langle #2|}
\providecommand{\Id}{\mathds{1}}
\providecommand{\one}{\mathds{1}}
\providecommand{\C}{\mathbb{C}}
\providecommand{\MN}[1]{M_{#1}(\C)}
\providecommand{\MD}{M_D(\C)}

% Verdict marker: \gapnote{<kind>}{<status>}. Kind is the policy
% classification (clarification, local-correction, scope-restriction,
% unfaithful, false-source, open-gap); status is open, wip (elimination
% underway), resolved, or historical. Machine-read by the site index;
% prints a verdict line.
\newcommand{\gapnote}[2]{\par\noindent\textbf{Verdict:} #1 (#2).\par}

\title{Full-Support Classification of Finite Dual Flux Paths}
\author{}
\date{}
\begin{document}
\maketitle
\gapnote{scope-restriction}{open}

Schuch, Cirac, and P\'erez-Garc\'ia,
\href{https://arxiv.org/abs/1001.3807v3}{arXiv:1001.3807v3},
Lemma 6.14, source \path{Papers/1001.3807/paper_v3.tex}, lines 2199--2214,
deforms a flux string by virtual invariance with its endpoints fixed.
On a torus, the topology of the allowed deformations must be specified.

Let $X=\mathbb Z_w\times\mathbb Z_h$, retain the four labelled unit steps,
and let $d(p)\in\mathbb Z^2$ be the sum of their signed displacements.
Let $p\sim_Xq$ be the relation generated by square moves and backtrack
cancellation, with every primal vertex allowed to be swept. For finite paths
with common endpoints,
\begin{align}
    p\sim_Xq\quad\Longleftrightarrow\quad d(p)=d(q).
    \label{eq:full_support_classification}
\end{align}
The forward direction follows by induction on the moves. For the converse,
choose $N=|p|+|q|$, place their common lifted start at $(N,N)$, and lift both
paths into $[0,2N]^2$. A prefix of length $k$ has coordinate displacement
bounded by $k$, so all lifted steps are defined in this rectangle. Equality
of total displacements identifies the lifted endpoints. The rectangular
comb normalizer gives the generated homotopy, without an embedding bound
on the rectangle. All four step labels are retained, including at torus
periods one and two. No winding sectors are identified merely because
their periodic endpoints agree.

The classification is formalized as
\leanid{TNLean.PEPS.torusDualHomotopy_univ_iff_displacement_eq}.
With virtual invariance at every site, actual physical state equality for
equal displacements follows as
\leanid{TNLean.PEPS.torusDualFluxState_eq_of_displacement_eq}.

The condition that every swept vertex is allowed is essential to the scope
of (\ref{eq:full_support_classification}). For a restricted swept set, holes
can obstruct a deformation. Neither completeness for restricted supports
nor construction of a path avoiding an arbitrary observation region is
asserted. Equality with a vacuum reduced density matrix requires a separate
argument, including suitable topology and boundary conditions: an annular
observation can surround an unobserved flux endpoint. These obligations
remain under \ghissue{8703}. Full-lattice Hamiltonian placement and its
common kernel remain separate under \ghissue{8947}.
\end{document}
